\section{Анализ предметной области и теоретические основы}

\subsection{Рынок недвижимости как объект компьютерного анализа}

Рынок жилой недвижимости относится к сферам, где одновременно действуют экономические закономерности, пространственная неоднородность и большой объём разрозненной информации. Для покупателя важны цена, район и состояние дома; для исследователя --- структура предложения, динамика цен, связь характеристик объекта с его стоимостью. Традиционные источники --- отчёты Росреестра, выборочные обследования, публикации аналитических агентств --- дают агрегированную картину, но редко позволяют оперативно построить произвольный срез по микропризнакам: этажности, типу дома, году постройки, удалённости от центра.

С распространением интернет-площадок ситуация изменилась. Агрегаторы объявлений ежедневно публикуют тысячи предложений о продаже квартир, и каждое объявление содержит набор атрибутов: стоимость, площадь, число комнат, адрес или координаты, параметры здания. Как показывают исследования, компьютерный анализ таких массивов помогает выявлять факторы ценообразования и особенности локальных сегментов рынка~\cite{soboleva2023,belyaev2015}. При этом сами площадки ориентированы на поиск жилья, а не на статистику: встроенные фильтры ограничены, сравнение нескольких сегментов рынка неудобно, географическая картина представлена фрагментарно.

Отсюда следует постановка задачи, характерная для прикладной информатики: необходимо не только получить данные из веб-источника, но и привести их к виду, пригодному для анализа, а затем обеспечить наглядное и гибкое исследование --- через специализированные программные средства визуализации и фильтрации~\cite{filatov2013,ahtimirov2024}. Далее в разделе рассматриваются теоретические аспекты каждого из этих этапов: от способов сбора информации в сети до принципов построения интерактивных аналитических панелей.

\subsection{Классификация источников данных о недвижимости}

С точки зрения анализа данных источники можно разделить на первичные и вторичные. К первичным относятся сведения, полученные непосредственно с площадок размещения объявлений или через обращение к ним программными средствами. К вторичным --- уже обработанные датасеты, отчёты и индексы, подготовленные третьими сторонами.

\textbf{Официальная статистика и реестры} дают высокую достоверность по сделкам и правам, но обновляются с задержкой и не всегда содержат детальные характеристики объекта в момент выставления на продажу.

\textbf{Коммерческие агрегаторы} (классифайды, порталы недвижимости) отражают текущее предложение: актуальные цены «запроса», а не факт сделки. Данные богаты атрибутами, но могут содержать дубликаты, устаревшие объявления и ошибки ввода со стороны продавца или агента.

\textbf{Открытые программные интерфейсы (API)} встречаются у крупных сервисов, однако для многих агрегаторов публичный API либо отсутствует, либо ограничен партнёрами. В таких случаях исследователь вынужден обращаться к методам извлечения данных из пользовательского веб-интерфейса.

\textbf{Пользовательские выгрузки} --- CSV, Excel, выгрузки из CRM агентств --- удобны для локального анализа, но плохо масштабируются и не гарантируют единообразие схемы полей.

Для задач мониторинга рынка в режиме, близком к реальному времени, наибольший интерес представляют агрегаторы объявлений. Их данные носят полуструктурированный характер: часть полей формализована (цена, площадь), часть --- текстовое описание (ремонт, инфраструктура), часть --- географические координаты или адрес в свободной форме. Это накладывает требования на этапы очистки, нормализации и контроля качества, рассмотренные в подразделе~1.4.

\subsection{Обзор методов сбора и анализа данных в веб-пространстве}

Сбор данных из интернета --- самостоятельное направление, на стыке программирования, сетевых технологий и прикладной статистики. В обзорной литературе и учебных курсах по анализу данных экосистема Python рассматривается как стандартный инструмент для последующей обработки табличных массивов~\cite{mckinney2023,vanderplas2023,coursera-python}. Библиотека pandas обеспечивает представление данных в виде таблицы (DataFrame), операции фильтрации, группировки и агрегирования; NumPy --- эффективные числовые вычисления над столбцами. Однако до этапа загрузки в DataFrame необходимо решить задачу \emph{извлечения} информации из HTML-страниц или API.

Веб-страница с точки зрения программы --- документ разметки HTML, дополненный каскадными таблицами стилей и скриптами JavaScript. Совокупность элементов страницы образует объектную модель документа (DOM). Аналитически полезные сведения могут находиться в текстовых узлах, атрибутах тегов, встроенных JSON-блоках (в том числе разметке schema.org) или подгружаться асинхронно после действий пользователя.

Различают два принципиально разных подхода к получению содержимого.

\textbf{Статический доступ} предполагает, что полный HTML уже содержит искомые данные. Достаточно выполнить HTTP-запрос, сохранить ответ и разобрать разметку парсером. Классическим инструментом для навигации по дереву тегов служит библиотека Beautiful Soup~\cite{beautifulsoup}: она позволяет находить элементы по имени, классу, идентификатору и извлекать текст или атрибуты.

\textbf{Динамический доступ} необходим, когда содержимое формируется на стороне клиента. Типичны одностраничные приложения (SPA): первичный HTML --- «оболочка», а объявления подгружаются запросами к внутренним API и вставляются в DOM скриптами. В такой ситуации простой HTTP-клиент получает пустой или неполный документ. Требуются средства автоматизации браузера --- в первую очередь Selenium WebDriver~\cite{selenium}, управляющий реальным или headless-браузером: открытие URL, прокрутка, клики, ожидание появления элементов, чтение DOM после выполнения JavaScript.

Сравнение подходов сведено в таблице~\ref{tab:parse-methods}.

\begin{table}[htbp]
\centering
\begin{tabular}{|p{3.2cm}|p{4.8cm}|p{4.8cm}|}
\hline
\textbf{Критерий} & \textbf{Статический парсинг} & \textbf{Автоматизация браузера} \\
\hline
Тип сайта & Серверный рендеринг HTML & SPA, AJAX, бесконечная прокрутка \\
\hline
Скорость & Выше & Ниже (накладные расходы браузера) \\
\hline
Устойчивость к JS & Нет & Да \\
\hline
Ресурсоёмкость & Низкая & Высокая \\
\hline
Типичные инструменты & HTTP-клиент, Beautiful Soup & Selenium, Playwright \\
\hline
\end{tabular}
\caption{Сравнение подходов к извлечению данных из веб-страниц}
\label{tab:parse-methods}
\end{table}

Для агрегаторов недвижимости всё чаще актуален второй подход или комбинация: браузер обеспечивает навигацию и подгрузку, после чего данные считываются из DOM или из перехваченных JSON-ответов.

После сбора начинается этап анализа. В учебниках по Data Science выделяют разведочный анализ данных (Exploratory Data Analysis, EDA) --- систематическое изучение распределений, пропусков, выбросов и связей между переменными до построения прогнозных моделей~\cite{demidova2022,grus2021}. Для описательной аналитики рынка недвижимости EDA нередко является конечной целью: исследователю нужны гистограммы цен, карта предложений, сравнение сегментов, а не обязательно модель машинного обучения. Вопросы интерпретируемости и участия эксперта в аналитическом процессе подробно обсуждаются в литературе по машинному обучению с участием человека~\cite{monarch2022}.

Для регулярного мониторинга рынка и совместного доступа к одним и тем же данным разовых скриптов недостаточно: целесообразно веб-приложение с интерактивной панелью~\cite{filimonova2024,sultanov2022}.

\subsection{Технологии парсинга и ограничения веб-источников}

\subsubsection{Структура HTML и локализация данных}

Проектирование парсера начинается с анализа структуры целевых страниц. HTML представляет иерархию вложенных элементов; CSS-классы и XPath-выражения служат для однозначного указания нужного фрагмента. На странице списка объявлений обычно повторяется один и тот же шаблон карточки; на странице объ объявления --- блоки «основные параметры», «о доме», «расположение». Разработчик фиксирует соответствие между полями будущей таблицы и фрагментами DOM.

Помимо видимого HTML данные могут дублироваться во встроенных JSON-LD-блоках (формат linked data для поисковых систем). Использование нескольких источников внутри одной страницы повышает полноту извлечения: если текстовое поле пусто, значение может быть получено из структурированной разметки.

\subsubsection{Поэтапный pipeline сбора}

На практике сбор с крупных каталогов редко сводится к одному проходу по сайту. Логично выделяют стадии:
\begin{enumerate}
    \item \textbf{сбор ссылок} --- обход страниц поиска с разными параметрами (район, число комнат), накопление URL карточек без детального разбора;
    \item \textbf{сбор деталей} --- последовательный обход списка ссылок, извлечение полного набора атрибутов каждого объявления;
    \item \textbf{нормализация} --- приведение «сырой» таблицы к единой схеме, расчёт производных показателей, удаление дубликатов;
    \item \textbf{контроль качества} --- проверка полноты ключевых полей, доли пропусков, географической правдоподобности координат.
\end{enumerate}
Такое разделение упрощает возобновление после сбоя, повторный прогон только одной стадии и параллельную настройку параметров обхода для разных регионов.

\subsubsection{Защита от автоматизированного доступа}

Владельцы коммерческих площадок ограничивают интенсивность автоматических обращений. К распространённым мерам относятся:
\begin{itemize}
    \item \textbf{ограничение частоты запросов} --- временные блокировки IP при превышении порога;
    \item \textbf{CAPTCHA} --- задачи, требующие участия человека;
    \item \textbf{проверка заголовков и отпечатка браузера} --- отличие скрипта от обычного пользователя;
    \item \textbf{динамическая структура разметки} --- изменение классов и путей DOM, затрудняющее поддержку парсера.
\end{itemize}
С теоретической точки зрения парсинг превращается в задачу устойчивого извлечения при неполной информации о внутреннем устройстве сайта и при активном противодействии~\cite{nesterov2024}. Рекомендуется варьировать паузы между запросами, имитировать естественную прокрутку, корректно обрабатывать сценарии с CAPTCHA (пауза, ручное решение, продолжение). Эти меры не гарантируют бессрочной работоспособности парсера, но снижают риск блокировки.

\subsubsection{Альтернатива парсингу --- программный интерфейс}

Если площадка предоставляет документированный API с авторизацией, предпочтительно использовать его: данные приходят в JSON или XML, схема полей известна, соблюдаются лимиты использования. Отсутствие открытого API вынуждает прибегать к парсингу пользовательского интерфейса, что повышает стоимость сопровождения при каждом изменении вёрстки сайта.

\subsection{Методология разведочного анализа и предобработки данных}

Сырые данные с веб-источника редко пригодны для немедленной визуализации. Ошибки ввода, пропуски, несогласованные форматы чисел и дат, текстовые поля вместо категорий --- типичные проблемы~\cite{mckinney2023}. Предобработка и EDA направлены на превращение такого массива в «чистую» таблицу с понятной семантикой столбцов.

\subsubsection{Пропущенные значения}

Механизмы появления пропусков классифицируют как полностью случайные (MCAR), случайные (MAR) и неслучайные (MNAR)~\cite{demidova2022}. Выбор стратегии зависит от доли пропусков и от того, связана ли их природа с целевым показателем (например, с ценой).

Возможные действия:
\begin{itemize}
    \item \textbf{удаление строк} --- оправдано при малой доле потерь в некритичных полях или когда пропуск означает отсутствие объекта на рынке;
    \item \textbf{удаление столбцов} --- если признак заполнен менее чем на заданный процент;
    \item \textbf{импутация} --- заполнение средним, медианой, модой или константой «не указано» для категорий;
    \item \textbf{маркировка} --- отдельный индикатор «значение отсутствовало», если пропуск сам несёт информацию.
\end{itemize}
Для цены и площади квартиры удаление неполных записей часто предпочтительнее заполнения средним, чтобы не искажать распределение.

\subsubsection{Выбросы и аномалии}

Выброс --- наблюдение, существенно отклоняющееся от основной массы данных. На рынке недвижимости выбросами могут быть опечатки (лишний ноль в цене), тестовые объявления или элитные объекты, не представительные для массового сегмента.

Классический метод --- межквартильный размах (IQR). Для переменной $X$ вычисляют первый $Q_1$ и третий $Q_3$ квартили, $IQR = Q_3 - Q_1$. Подозрительными считают значения вне интервала
\begin{equation}
    [Q_1 - 1{,}5 \cdot IQR;\; Q_3 + 1{,}5 \cdot IQR].
    \label{eq:iqr}
\end{equation}
Альтернатива --- z-score: отклонение более чем на $3\sigma$ от среднего. Медиана и квартили устойчивее к выбросам, чем среднее арифметическое~\cite{vanderplas2023,grus2021}.

Решение об удалении или сохранении выброса должно опираться на предметную область: уникальный пентхаус и ошибка парсера требуют разной реакции.

\subsubsection{Распределения и описательная статистика}

Стоимость жилья на вторичном рынке часто имеет правостороннее асимметричное распределение: много объектов в «среднем» сегменте и длинный «хвост» дорогих предложений. Логнормальная модель нередко лучше описывает цены, чем нормальная. Для описания центра распределения в таких случаях предпочитают \textbf{медиану}; среднее арифметическое чувствительно к дорогим объектам. \textbf{Мода} полезна для дискретных признаков (число комнат).

Понимание формы распределения влияет на выбор визуализации (шаг гистограммы, шкала осей) и на интерпретацию трендов.

\subsubsection{Корреляция и связи между признаками}

Между площадью и ценой часто наблюдается положительная связь. Коэффициент корреляции Пирсона $r$ измеряет линейную зависимость двух количественных переменных:
\begin{equation}
    r = \frac{\sum (x_i - \bar{x})(y_i - \bar{y})}{\sqrt{\sum (x_i - \bar{x})^2 \sum (y_i - \bar{y})^2}}.
    \label{eq:pearson}
\end{equation}
Значения $r$ близкие к $+1$ или $-1$ указывают на сильную линейную связь; близкие к нулю --- на её отсутствие. Корреляция не доказывает причинно-следственную связь и плохо улавливает нелинейные зависимости, поэтому её дополняют диаграммами рассеяния~\cite{monarch2022}.

\subsubsection{Инженерия признаков}

Исходные поля объявления можно обогатить производными признаками, облегчающими анализ:
\begin{itemize}
    \item \textbf{удельная цена} --- отношение стоимости к площади (руб./м²);
    \item \textbf{возраст здания} --- разность между текущим годом и годом постройки;
    \item \textbf{относительный этаж} --- отношение этажа к этажности дома;
    \item \textbf{расстояние до условного центра города} --- по формуле haversine из географических координат;
    \item \textbf{бинарные индикаторы} --- наличие балкона, лифта, парковки и т.п.
\end{itemize}

Категориальные переменные (тип дома: панельный, кирпичный, монолитный) для алгоритмов фильтрации и некоторых моделей преобразуют в числовой вид. \textbf{One-hot encoding} создаёт отдельный столбец для каждой категории со значениями 0 или 1. Альтернатива --- порядковое кодирование, если категории имеют естественный порядок (например, состояние ремонта).

\subsubsection{Дедупликация и слияние записей}

Одно и то же объявление может быть собрано повторно при разных запусках парсера. Дубликаты искажают частоты и средние показатели. Идентификация ведётся по стабильному ключу --- URL объявления, внешнему идентификатору на площадке или комбинации адреса, площади и цены. При слиянии двух версий одной записи из разных выгрузок логично объединять непустые поля: если в первой версии не было координат, а во второй они появились, результирующая строка должна их сохранить.

\subsubsection{Контроль качества данных}

Качество данных оценивают по нескольким измерениям~\cite{demidova2022}:
\begin{itemize}
    \item \textbf{полнота} --- доля заполненных значений в ключевых столбцах;
    \item \textbf{уникальность} --- отсутствие нежелательных дубликатов;
    \item \textbf{валидность} --- соответствие форматов и допустимых диапазонов (цена ${>}0$, площадь в разумных пределах);
    \item \textbf{согласованность} --- непротиворечивость связанных полей (этаж не больше этажности дома);
    \item \textbf{актуальность} --- давность сбора относительно момента анализа.
\end{itemize}
Формализованные пороги (например, не менее 75\% заполненности координат, не более 15\% пропусков в цене) позволяют автоматически отклонять низкокачественные выгрузки до этапа визуализации.

\subsection{Теоретические основы визуализации данных}

Визуализация переводит числовую таблицу в форму, удобную для зрительного восприятия. Человек быстрее замечает тренды, кластеры и аномалии на графике, чем при просмотре тысяч строк~\cite{vanderplas2023}. Цель визуализации в аналитике недвижимости --- не украшение отчёта, а снижение когнитивной нагрузки при исследовании многомерных данных.

\subsubsection{Принципы эффективного графического представления}

При выборе типа диаграммы учитывают вид переменных:
\begin{itemize}
    \item одна количественная --- гистограмма, box plot, плотность;
    \item одна категориальная --- столбчатая диаграмма частот;
    \item две количественные --- диаграмма рассеяния;
    \item количественная и категориальная --- grouped bar chart, box plot по группам;
    \item географические координаты --- точечная карта, тепловая карта плотности.
\end{itemize}
Нарушение соответствия типа графика типу данных (например, гистограмма для номинальных категорий без порядка) ведёт к неверной интерпретации.

\subsubsection{Гистограммы и биннинг}

Гистограмма делит диапазон значений на интервалы (bins) и показывает число наблюдений в каждом. Ширина интервала влияет на вид графика: слишком узкие bins дают шум, слишком широкие скрывают мультимодальность.

Правило Стерджеса предлагает число интервалов $k = \lceil \log_2 n \rceil + 1$ для выборки из $n$ точек. Правило Фридмана--Диакониса учитывает размах данных и межквартильный размах:
\begin{equation}
    h = 2 \cdot IQR \cdot n^{-1/3},
    \label{eq:fd}
\end{equation}
где $h$ --- ширина bin~\cite{grus2021}. На практике допустима интерактивная подстройка шага пользователем.

\subsubsection{Диаграммы рассеяния и проблема наслоения}

При построении scatter-графика «цена--площадь» для большой выборки точки накладываются друг на друга (overplotting), и плотность участков становится неразличимой. Приёмы борьбы: полупрозрачность маркеров (alpha blending), субдискретизация (random sampling), контурные линии плотности, шестигранная бининг-диаграмма (hexbin)~\cite{mckinney2023}.

Добавление линии тренда (например, методом наименьших квадратов) облегчает оценку общего направления связи, но не должно подменять анализ остатков и выбросов.

\subsubsection{Кодирование категорий цветом и формой}

Для различения групп (тип дома, район) используются визуальные каналы: цвет, форма маркера, тип линии. Рекомендуется не перегружать один график большим числом категорий --- зрительное восприятие уверенно различает ограниченное количество оттенков. Категориальные шкалы предпочтительнее раскрашивать качественными палитрами без ложного порядка «больше-меньше»~\cite{monarch2022}.

\subsection{Интерактивные аналитические панели (дашборды)}

\subsubsection{Отличие дашборда от статического отчёта}

Статический отчёт (PDF, слайды, фиксированная таблица Excel) фиксирует один срез данных на момент подготовки. \textbf{Интерактивный дашборд} позволяет пользователю менять фильтры, масштаб и состав виджетов, получая обновлённую картину без повторного построения документа~\cite{filatov2013,filimonova2024}. Связь «действие пользователя $\rightarrow$ пересчёт выборки $\rightarrow$ обновление графиков» лежит в основе современных систем бизнес-аналитики и специализированных исследовательских интерфейсов.

Преимущества интерактивности для анализа недвижимости:
\begin{itemize}
    \item быстрая проверка гипотез («как меняется медиана цены при ограничении радиуса от центра»);
    \item сопоставление сегментов («однокомнатные vs двухкомнатные») на одном экране;
    \item географическая привязка --- карта как дополнение к табличным и графическим представлениям;
    \item доступность для пользователей без навыков программирования.
\end{itemize}

\subsubsection{Архитектура представления данных в дашборде}

Концептуально дашборд опирается на три слоя:
\begin{enumerate}
    \item \textbf{хранилище} --- таблица наблюдений (строки --- объявления, столбцы --- признаки);
    \item \textbf{логика фильтрации} --- правила отбора строк по диапазонам и категориям;
    \item \textbf{слой представления} --- виджеты (KPI, графики, карта, таблица), каждый из которых получает уже отфильтрованную или агрегированную выборку.
\end{enumerate}
Изменение фильтра пересчитывает все зависимые виджеты. При большом объёме данных критична производительность: отбор строк и агрегирование должны выполняться за время, не нарушающее плавность интерфейса (ориентир --- доли секунды на типичной выборке).

\subsubsection{Типовые элементы дашборда для рынка недвижимости}

\textbf{KPI-блок (ключевые показатели)} --- краткие числа по текущей выборке: число объектов, медиана цены, медиана цены за м², медиана площади. Медианы предпочтительнее средних при асимметричных распределениях цен.

\textbf{Панель фильтров} --- элементы ввода для цены, площади, комнатности, года постройки, типа дома, расстояния от центра. Дополнительные условия по произвольным столбцам расширяют гибкость без изменения кода приложения.

\textbf{Набор графиков} --- распределения, сравнения категорий, scatter с трендом. Возможность включать и отключать отдельные графики снижает перегруженность экрана.

\textbf{Карта} --- точечное отображение объектов с кластеризацией при большой плотности; всплывающие подсказки с ценой и адресом.

\textbf{Табличное представление} --- детальный просмотр отдельных объявлений с сортировкой по столбцам и ссылкой на первоисточник.

\textbf{Сравнение сегментов} --- два независимых набора фильтров и параллельное отображение KPI или графиков для каждого.

\subsubsection{Географическая визуализация}

Недвижимость по своей природе связана с географией. Карта дополняет табличный анализ: видны «горячие» зоны предложения, радиальные градиенты цены, изолированные кластеры. Точечные карты подходят для отдельных объявлений; хороплетные карты (раскраска районов по агрегированному показателю) --- для сравнения территорий, если выборка достаточна для устойчивых средних.

При отсутствии текстового адреса, но наличии координат применяют \textbf{обратное геокодирование} (reverse geocoding) --- преобразование пары (широта, долгота) в адрес или название места с помощью внешних геосервисов. Результаты кэшируют, чтобы не обращаться к API повторно для одних и тех же точек.

\subsubsection{Принципы проектирования интерфейса}

При проектировании аналитической панели рекомендуется:
\begin{itemize}
    \item отделить область управления (фильтры) от области результатов (графики, карта);
    \item обеспечить немедленную обратную связь при изменении параметров;
    \item поддерживать адаптивную вёрстку для разных размеров экрана;
    \item предусмотреть светлую и тёмную тему для длительной работы;
    \item не перегружать один экран: при необходимости использовать вкладки или сворачиваемые блоки.
\end{itemize}
Подход, при котором основной экран сразу показывает сводку и виджеты («dashboard first»), распространён в системах мониторинга и исследовательских UI~\cite{filimonova2024}.

\subsection{Обзор программных средств анализа данных}

\subsubsection{Язык Python и библиотеки обработки таблиц}

Python закрепился в качестве универсального языка для анализа данных благодаря библиотекам общего назначения~\cite{habr-bigdata,w3schools-numpy}. Pandas предоставляет структуры DataFrame и Series, операции join, groupby, pivot; NumPy --- векторизованные вычисления. Для задач недвижимости типичны операции: группировка по району, расчёт медианы цены за м², фильтрация по квантилям.

\subsubsection{Среды разработки программных комплексов}

Проект, объединяющий Python-парсер, серверную часть и веб-клиент, удобно вести в интегрированной среде разработки (IDE), поддерживающей несколько языков и запуск скриптов из терминала. Такие среды (Visual Studio Code, JetBrains WebStorm и аналоги) дают подсветку синтаксиса, навигацию по проекту, отладку и единый каталог для всех компонентов системы. Для многопользовательской аналитики и разграничения прав всё равно требуется полноценное веб-приложение, а не локальные скрипты~\cite{ahtimirov2024}.

\subsubsection{Веб-технологии для интерактивной визуализации}

Клиентские библиотеки (D3.js, Chart.js, Recharts и др.) строят графики в браузере, реагируя на изменение данных. Серверная часть может ограничиваться выдачей JSON и авторизацией; фильтрация и отрисовка выполняются на клиенте --- паттерн, характерный для одностраничных приложений~\cite{sultanov2022,nesterov2024}. Выбор конкретного фреймворка определяется требованиями к типизации, экосистеме компонентов и производительности; теоретически любая declarative charting library, интегрируемая с реактивной моделью UI, пригодна для дашборда.

\subsection{Выводы по первому разделу}

В первом разделе рассмотрены теоретические основы, необходимые для построения системы анализа рынка недвижимости на основе веб-данных:
\begin{enumerate}
    \item агрегаторы объявлений --- актуальный, но «шумный» источник; данные требуют очистки и контроля качества;
    \item для динамических сайтов обоснован переход от статического парсинга HTML к автоматизации браузера; сбор целесообразно организовать поэтапно;
    \item предобработка опирается на методы описательной статистики: работа с пропусками, выбросами, производными признаками, кодирование категорий;
    \item интерактивный дашборд даёт гибкость, недостижимую для статического отчёта; его элементы (KPI, фильтры, графики, карта) согласованы с типами анализируемых признаков;
    \item программная реализация может опираться на связку Python (сбор и ETL) и веб-стек (визуализация и доступ пользователей).
\end{enumerate}
